% Full proof attachment. Owned by counterfactual_algebra. Labels: MFC.
% Original sources: mixed_support_ledger.tex (ML), lines 182--338;
% split_support_absolute_arithmetic_curve_v5.tex (V5), lines 591--659
% and 1316--1556. All carriers and operations below are explicit.
\providecommand{\Chi}{\mathrm{X}}
\section{Full mixed carriers, their faces, prime branches, and depth maps}
\label{sec:MFC}

\paragraph{MFC.1. Conventions and the original split scalar.}
All rings in this attachment are commutative with identity and are nonzero.
All semiring homomorphisms preserve the specified zero and unit unless a
different category is explicitly named. A semiring ideal contains zero, is
closed under addition, and absorbs multiplication by arbitrary elements;
no subtractivity requirement is imposed. A prime ideal is a proper ideal
whose complement is multiplicatively closed. The length of a strict prime
chain is its number of strict inclusions.
For a ring \(R\), write
\[
 G(R)=\{\tau_R\}\sqcup\{r^\bullet:r\in R\},\qquad
 e_R=0_R^\bullet,\qquad 1_{G(R)}=1_R^\bullet.                 \tag{MFC.1}
\]
The union is disjoint. Its operations are
\[
 r^\bullet+s^\bullet=(r+s)^\bullet,\quad
 r^\bullet s^\bullet=(rs)^\bullet,\quad
 \tau_R+x=x,\quad \tau_Rx=\tau_R.
                                                               \tag{MFC.2}
\]
In particular, a cancellation of supported amplitudes produces \(e_R\),
not \(\tau_R\).
These operations make a commutative semiring. On supported entries the
laws are the ring laws. In a sum an occurrence of \(\tau_R\) can be deleted,
and a product with an occurrence of \(\tau_R\) equals \(\tau_R\); these
observations prove both associativities in every other case. For
distributivity the new cases are
\(\tau_R(x+y)=\tau_R=\tau_Rx+\tau_Ry\) and
\(x(\tau_R+y)=xy=x\tau_R+xy\).
The amplitude and support maps
\[
 \begin{aligned}
 p_R:G(R)&\longrightarrow R,&
 p_R(\tau_R)&=0_R,& p_R(r^\bullet)&=r,\\
 \chi_R:G(R)&\longrightarrow\mathbb B,&
 \chi_R(\tau_R)&=0,&\chi_R(r^\bullet)&=1
 \end{aligned}                                           \tag{MFC.3}
\]
are unital semiring homomorphisms; here
\(\mathbb B=\{0,1\}\) has join as addition and meet as multiplication.
Indeed (MFC.2) proves the amplitude identities, and a sum is absent
exactly when both terms are absent, whereas a product is absent exactly
when at least one factor is absent.

\paragraph{MFC.2. Independent Cartesian faces and their entire fibres.}
Let \(I\) be a nonempty finite set and \(R_i\) a nonzero ring for each
\(i\in I\). Put
\[
 P_I=\prod_{i\in I}G(R_i),\qquad R_I=\prod_{i\in I}R_i,\qquad
 p=\prod_i p_{R_i},\qquad
 s(x)=\{i:x_i\ne\tau_{R_i}\}.                              \tag{MFC.4}
\]
Both operations of \(P_I\) are coordinatewise. The map
\(s:P_I\to\mathcal P(I)\) is a unital semiring homomorphism for
union and intersection on \(\mathcal P(I)\), by the last two assertions
of MFC.1 in each coordinate. Its exact face over \(A\subseteq I\) is
\[
 s^{-1}(A)=
 \prod_{i\in A}\{r_i^\bullet:r_i\in R_i\}
 \times\prod_{i\notin A}\{\tau_{R_i}\}.                      \tag{MFC.5}
\]
The ordering of factors is the original ordering indexed by \(I\);
the displayed expression specifies each coordinate rather than changing
that ordering.
For \(r=(r_i)\in R_I\), put \(J(r)=\{i:r_i\ne0_{R_i}\}\). For each
\(A\supseteq J(r)\) define
\[
 x(r,A)_i=
 \begin{cases}r_i^\bullet&i\in A,\\
               \tau_{R_i}&i\notin A.
 \end{cases}
 \qquad
 p^{-1}(r)=\{x(r,A):J(r)\subseteq A\subseteq I\}.             \tag{MFC.6}
\]
To prove equality, an absent coordinate of any preimage must have zero
amplitude; hence its mask contains \(J(r)\). At an active coordinate its
amplitude forces the unique supported entry \(r_i^\bullet\).
Conversely each displayed tuple has amplitude \(r\). Distinct masks give
distinct tuples, including at zero amplitude. Consequently
\[
 |p^{-1}(r)|=2^{|I|-|J(r)|},\qquad
 (p,s):P_I\hookrightarrow R_I\times\mathcal P(I)
 \text{ has image }\{(r,A):J(r)\subseteq A\}.                \tag{MFC.7}
\]
In particular the amplitude-zero fibre has all \(2^{|I|}\) faces.
These statements retain supported zeros on every proper face.

\paragraph{MFC.3. Synchronized inclusion and the precise failed collapse.}
The original synchronized inclusion is
\[
 j:G(R_I)\hookrightarrow P_I,\qquad
 j(\tau_{R_I})=(\tau_{R_i})_i,\qquad
 j(r^\bullet)=(r_i^\bullet)_i.                             \tag{MFC.8}
\]
It is injective: the absent image differs from every fully supported
image, and amplitudes distinguish fully supported images. Applying
(MFC.2) coordinatewise proves preservation of sum, product, zero, and
unit. Its image is exactly \(s^{-1}(\{\varnothing,I\})\).
Thus the square
\[
 \begin{array}{ccc}
 G(R_I)&\xrightarrow{\ j\ }&P_I\\
 {\scriptstyle\chi_{R_I}}\downarrow&&\downarrow{\scriptstyle s}\\
 \mathbb B&\xrightarrow{\ \delta\ }&\mathcal P(I)
 \end{array}
 \qquad \delta(0)=\varnothing,\quad\delta(1)=I              \tag{MFC.9}
\]
is a pullback of semirings. Explicitly, an element of the fibre product
is a pair \((\epsilon,x)\) with mask empty when \(\epsilon=0\) and full
when \(\epsilon=1\). In the first case it comes uniquely from
\(\tau_{R_I}\); in the second case its amplitudes give its unique
preimage \(r^\bullet\). This bijection respects all operations by
(MFC.8), proving the asserted universal property. Also \(p j=p_{R_I}\).

For \(I=\{1,2\}\), there is no semiring homomorphism
\[
 q:P_I\longrightarrow G(R_1\times R_2),\qquad
 p_{R_1\times R_2}q=p,
 \quad q(x)\text{ supported whenever }s(x)\ne\varnothing.
                                                               \tag{MFC.10}
\]
Indeed let \(x=(1_{R_1}^\bullet,\tau_{R_2})\) and
\(y=(\tau_{R_1},1_{R_2}^\bullet)\). Their product is the specified zero
of \(P_I\). The stipulated images are \((1,0)^\bullet\) and
\((0,1)^\bullet\), whose product is \(e_{R_1\times R_2}\), distinct
from \(q(0_{P_I})=\tau_{R_1\times R_2}\). This contradicts
multiplicativity and zero preservation. The statement concerns exactly
the map in (MFC.10); it leaves the proven maps \(j,p,s\) intact.

\paragraph{MFC.4. The original quadratic convolution carrier.}
Fix a nonzero ring \(A\), set \(S=G(A)\),
\(u=(-1_A)^\bullet\), and put
\[
 D_A=S\oplus S,\qquad
 (a,b)\star(c,d)=(ac+ubd,\ ad+bc),\qquad
 \mathbf0=(\tau_A,\tau_A),\quad
 X_1=(1_A^\bullet,\tau_A).                                \tag{MFC.11}
\]
Addition remains coordinatewise. In this formula \(a,b,c,d\) lie in
\(S\), so every sum and product on the right is the original split
operation. Set
\[
 B=A[t]/(t^2+1),\qquad
 X_a=(a^\bullet,\tau_A),\quad
 Y_b=(\tau_A,b^\bullet),\quad
 Z_{a+tb}=(a^\bullet,b^\bullet).                            \tag{MFC.12}
\]
Division by the monic polynomial \(t^2+1\) gives every element of \(B\)
a unique expression \(a+tb\). For completeness, repeated subtraction
of the leading monomial times \(t^2+1\) reduces a polynomial to degree
at most one. A nonzero multiple of this monic polynomial has degree
at least two, since its leading coefficient equals the nonzero leading
coefficient of the multiplier. Therefore two reduced expressions cannot
differ by a nonzero multiple. In particular \(A\hookrightarrow B\),
\(B\ne0\), and \(t\) is a unit with inverse \(-t\). No assumption on
the characteristic or on the invertibility of \(2\) is used.

The multiplication in (MFC.11) is associative. For
\((a,b),(c,d),(h,k)\in S^2\), the first coordinate in either
association is
\[
 ach+ubdh+uadk+ubck,
 \qquad\text{and the second is}\qquad
 ack+ubdk+adh+bch.                                       \tag{MFC.13}
\]
This follows by distributing each of the two products; commutativity
of \(S\) puts the resulting four terms in the displayed order.
Commutativity follows directly from (MFC.11). Distributing its two
coordinates proves distributivity over coordinatewise addition.
Substitution gives the stated zero and unit, so \(D_A\) is a
commutative semiring. Substitution also gives every nonempty-face
product:
\[
 \begin{array}{lll}
 X_aX_c=X_{ac},&X_aY_b=Y_{ab},&Y_bY_d=X_{-bd},\\
 X_aZ_z=Z_{az},&Y_bZ_z=Z_{tbz},&Z_zZ_w=Z_{zw}.
 \end{array}                                             \tag{MFC.14}
\]
For example, for \(z=c+td\), \(Y_bZ_z\) has coordinates
\((-bd)^\bullet,(bc)^\bullet\), which are exactly those of \(tbz\).
Products of two full faces similarly have coordinates
\((ac-bd)^\bullet,(ad+bc)^\bullet\).
The four elements \(\mathbf0,X_0,Y_0,Z_0\) are distinct.
Unlike \(P_{\{1,2\}}\), the carrier \(D_A\) uses convolution
multiplication; (MFC.14) specifies the difference without identifying
the two semiring structures.

\paragraph{MFC.5. Four support fibres and the scalar character.}
Write \(\mathbb B[C_2]\) for the semiring of pairs of bits with
coordinatewise join and
\[
 (p,q)(r,s)=(pr\vee qs,\ ps\vee qr).
 \qquad
 \Chi:D_A\longrightarrow\mathbb B[C_2],\quad
 \Chi(a,b)=(\chi_A(a),\chi_A(b)).                          \tag{MFC.15}
\]
Here \(pr\) means Boolean meet. Applying \(\chi_A\) to (MFC.11),
and using \(\chi_A(u)=1\), proves the multiplication formula;
addition, zero, and unit are also preserved. The support semiring
laws can alternatively be verified by distributing Boolean products
as in (MFC.13). The fibres over \((0,0),(1,0),(0,1),(1,1)\)
are respectively \(\{\mathbf0\},\{X_a:a\in A\},
\{Y_b:b\in A\},\{Z_z:z\in B\}\). Every one retains its zero-amplitude
entry. The amplitude map
\(\pi:D_A\to B\), obtained by replacing absent coordinates by zero,
preserves all operations by (MFC.11). The map \((\pi,\Chi)\) is
injective: support determines which coordinates are absent, and the
unique coordinates of \(\pi(x)\) determine all active entries.

There is exactly one unital semiring homomorphism \(f:D_A\to\mathbb B\):
it sends \(\mathbf0\) to \(0\) and every nonempty face to \(1\).
To prove uniqueness, \(f(X_1)=1\) and
\(X_1+X_{-1}=X_0\) imply \(f(X_0)=1\).
The equality \(X_0X_a=X_0\) then forces \(f(X_a)=1\) for every \(a\).
Next \(Y_1^2=X_{-1}\) forces \(f(Y_1)=1\), and
\(Y_a=X_aY_1\) forces \(f(Y_a)=1\).
Finally \(Z_{a+tb}=X_a+Y_b\) forces \(f(Z_{a+tb})=1\).
Existence follows because a sum or product of two nonempty faces
is nonempty, as follows respectively from coordinate addition and
(MFC.14); the cases with \(\mathbf0\) are the zero laws.
Thus scalar support is an exact further observation of the four
fibres, and does not replace (MFC.15).

\paragraph{MFC.6. The quadratic synchronization quotient and all fibres.}
Define
\[
 q_D:D_A\longrightarrow G(B),\quad
 \mathbf0\longmapsto\tau_B,\quad X_a\longmapsto a^\bullet,\quad
 Y_b\longmapsto(tb)^\bullet,\quad Z_z\longmapsto z^\bullet.
                                                               \tag{MFC.16}
\]
This is a surjective unital semiring homomorphism. Surjectivity follows
from the \(Z\)-face and \(\mathbf0\). A sum of two nonempty faces is
nonempty and has amplitude equal to the sum in \(B\). The same assertion
for multiplication follows from the six products in (MFC.14).
Thus both operations are preserved, even when the resulting amplitude
is zero. Operations with \(\mathbf0\) obey the zero laws, and
\(q_D(X_1)=1_B^\bullet\).
Its complete fibre list, in the unique coordinates \(z=a+tb\), is
\[
 \begin{array}{c|c}
 \text{target}&\text{fibre}\\ \hline
 \tau_B&\{\mathbf0\}\\
 e_B&\{X_0,Y_0,Z_0\}\\
 (a+tb)^\bullet,\ a\ne0,\ b\ne0&\{Z_{a+tb}\}\\
 a^\bullet,\ a\ne0&\{X_a,Z_a\}\\
 (tb)^\bullet,\ b\ne0&\{Y_b,Z_{tb}\}.
 \end{array}                                             \tag{MFC.17}
\]
Indeed an \(X\)-preimage requires second coordinate zero, a
\(Y\)-preimage requires first coordinate zero, and the \(Z\)-preimage
always exists uniquely. These alternatives exhaust the four faces.

The kernel congruence of \(q_D\) is generated by \(X_0\sim Y_0\).
Adding \(X_0\) and then \(Y_0\) to this relation shows
\(X_0\sim Z_0\sim Y_0\). Adding \(X_a\) to \(X_0\sim Y_0\)
gives \(X_a\sim Z_a\); adding \(Y_b\) gives
\(Z_{tb}\sim Y_b\). Hence each fibre in (MFC.17) is identified.
Conversely the kernel of the proved homomorphism \(q_D\) is a
congruence containing the generator, so no generated identification
can connect distinct fibres. This proves the equality of congruences.

\paragraph{MFC.7. The full section, its unit, and the localization.}
There is a zero-preserving additive and multiplicative section
\[
 \iota:G(B)\longrightarrow D_A,\qquad
 \iota(\tau_B)=\mathbf0,\quad \iota(z^\bullet)=Z_z,\qquad
 q_D\iota=\operatorname{id}_{G(B)}.                        \tag{MFC.18}
\]
Addition within the \(Z\)-face is \(Z_z+Z_w=Z_{z+w}\);
multiplication is (MFC.14), and \(\mathbf0\) obeys the required
zero laws. These prove every asserted property. This section is
not unital as a map to \(D_A\), since
\(\iota(1_B^\bullet)=Z_1=(1_A^\bullet,0_A^\bullet)\ne X_1\).

In fact no unital semiring section \(\sigma\) of \(q_D\) exists.
The equality \(1_B^\bullet+e_B=1_B^\bullet\) forces
\(X_1+\sigma(e_B)=X_1\). Among the three possibilities in the
\(e_B\)-fibre, only \(X_0\) satisfies this equality: adding \(Y_0\)
or \(Z_0\) gives \(Z_1\). Thus \(\sigma(e_B)=X_0\).
The \(t^\bullet\)-fibre is \(\{Y_1,Z_t\}\), using \(1_A\ne0_A\).
But \(e_Bt^\bullet=e_B\), while
\(X_0Y_1=Y_0\ne X_0\) and \(X_0Z_t=Z_0\ne X_0\).
Both possibilities contradict multiplicativity.

The element \(Z_1\) is idempotent, and its corner is
\[
 Z_1D_A=\{\mathbf0\}\sqcup\{Z_z:z\in B\},\qquad
 D_A[Z_1^{-1}]\ \cong\ Z_1D_A\ \cong\ G(B).               \tag{MFC.19}
\]
The corner has unit \(Z_1\). The map \(d\mapsto Z_1d\)
is a unital homomorphism to that corner: idempotence gives
\((Z_1d)(Z_1d')=Z_1dd'\), and distributivity gives additivity.
It equals \(\iota q_D\) with this codomain, by (MFC.14).
To prove the localization universal property, let \(f:D_A\to T\)
be a unital homomorphism for which \(f(Z_1)\) is invertible.
An invertible idempotent is \(1_T\), since multiplying
\(f(Z_1)^2=f(Z_1)\) by its inverse gives that equality.
Therefore \(f(d)=f(Z_1d)\). Restricting \(f\) to the corner gives
the unique unital homomorphism through which \(f\) factors; it is
unital because \(f(Z_1)=1_T\), and uniqueness follows from
surjectivity of \(d\mapsto Z_1d\).
This proves (MFC.19), retaining both the full fibres of \(q_D\)
and the exact category of its section.

\paragraph{MFC.8. All ordinary semiring ideals of \(D_A\).}
For ring ideals \(H\triangleleft B\), \(J\triangleleft A\), define
\[
 \begin{aligned}
 M_H&=\{\mathbf0\}\cup\{Z_z:z\in H\},\qquad M=M_B,\\
 I_{J,H}&=M_H\cup\{X_a:a\in J\}\cup\{Y_a:a\in J\},
               \qquad JB\subseteq H .
 \end{aligned}                                           \tag{MFC.20}
\]
Here \(JB\) is the ideal of \(B\) generated by the image of \(J\).
Every ideal of \(D_A\) is exactly one of
\[
 \{\mathbf0\},\qquad M_H,\qquad I_{J,H}
 \quad(H\triangleleft B,\ J\triangleleft A,\ JB\subseteq H).
                                                               \tag{MFC.21}
\]
We prove the statement for the ordinary ideal convention in MFC.1.
Let \(\mathcal I\) be such an ideal. If it has no \(X\)- or \(Y\)-entry
and is not \(\{\mathbf0\}\), it contains some \(Z_z\). The set
\(H=\{z:Z_z\in\mathcal I\}\) is nonempty. Multiplication by
\(Z_0\) shows \(0\in H\), multiplication by \(Z_{-1}\) gives
additive inverses, addition gives closure under sums, and
multiplication by \(Z_w\) for arbitrary \(w\in B\) gives absorption.
Hence \(H\) is a ring ideal and \(\mathcal I=M_H\).

If \(\mathcal I\) has an \(X\)- or \(Y\)-entry, put
\(J_X=\{a:X_a\in\mathcal I\}\) and \(J_Y=\{a:Y_a\in\mathcal I\}\).
The identities \(X_aY_1=Y_a\) and \(Y_aY_{-1}=X_a\)
give \(J_X=J_Y=:J\), and this common set is nonempty.
Multiplication of an \(X\)-entry by \(X_0\) puts \(0\) in \(J\);
addition, multiplication by \(X_{-1}\), and multiplication by
arbitrary \(X_r\) prove respectively closure under sums, additive
inverses, and absorption by \(A\). Thus \(J\triangleleft A\).
Multiplying an \(X\)-entry by \(Z_0\) puts \(Z_0\) in
\(\mathcal I\), and the preceding \(Z\)-face argument proves that
\(H=\{z:Z_z\in\mathcal I\}\) is a ring ideal of \(B\).
For \(a\in J,w\in B\), \(X_aZ_w=Z_{aw}\) puts \(aw\) in \(H\);
closure under finite sums proves \(JB\subseteq H\).
The definitions now give \(\mathcal I=I_{J,H}\).

Conversely \(M_H\) is closed under addition, and every nonempty face
multiplied by \(Z_h\), \(h\in H\), remains a \(Z\)-entry with
amplitude in \(H\), by (MFC.14). The absent face multiplies
it to \(\mathbf0\), which is also in \(M_H\). Thus \(M_H\) is an ideal.
For \(I_{J,H}\), same-face sums stay in \(J\) or \(H\).
Cross-face sums have the precise forms
\[
 X_a+Y_b=Z_{a+tb},\qquad X_a+Z_h=Z_{a+h},\qquad
 Y_b+Z_h=Z_{tb+h}.                                       \tag{MFC.22}
\]
When \(a,b\in J\) and \(h\in H\), these amplitudes lie in \(H\)
because \(JB\subseteq H\). Multiplication of an \(X\)- or
\(Y\)-entry of \(I_{J,H}\) by an \(X\)- or \(Y\)-entry of \(D_A\)
has coefficient in \(J\), and its product with a \(Z\)-entry
has amplitude in \(JB\). A \(Z\)-entry of \(I_{J,H}\) multiplied
by any face remains in \(M_H\). Products with \(\mathbf0\)
are \(\mathbf0\). This verifies absorption in every case.
Each listed family is therefore an ideal. Their descriptions
are unique by reading the amplitudes in each face. The three
families are disjoint: \(M_H\) contains \(Z_0\) but no \(X\)- or
\(Y\)-entry, while \(I_{J,H}\) contains \(X_0,Y_0,Z_0\).

Their complete inclusion rules, aside from the least ideal
\(\{\mathbf0\}\), are
\[
 \begin{aligned}
 M_H\subseteq M_{H'}&\iff H\subseteq H',\\
 M_H\subseteq I_{J',H'}&\iff H\subseteq H',\\
 I_{J,H}\subseteq I_{J',H'}&\iff J\subseteq J'
                                      \text{ and }H\subseteq H'.
 \end{aligned}                                           \tag{MFC.23}
\]
No \(I_{J,H}\) is contained in an \(M_{H'}\), since it contains
\(X_0\). Each equivalence follows by inspecting the relevant
faces in (MFC.20), which are disjoint.

\paragraph{MFC.9. All prime branches and their exact inclusions.}
For \(\mathfrak q\in\operatorname{Spec}B\) and
\(\mathfrak p\in\operatorname{Spec}A\), put
\[
 P_\tau=\{\mathbf0\},\qquad
 Q_{\mathfrak q}=I_{\mathfrak q\cap A,\mathfrak q},
 \qquad T_{\mathfrak p}=I_{\mathfrak p,B}.
 \qquad
 \operatorname{Spec}D_A=
 \{P_\tau,M\}\sqcup\{Q_{\mathfrak q}\}\sqcup\{T_{\mathfrak p}\}.
                                                               \tag{MFC.24}
\]
We prove exhaustion and primality, rather than imposing a subtractive
prime convention. The complement of \(P_\tau\) is closed under
multiplication by (MFC.14). The complement of \(M\) consists
of all \(X\)- and \(Y\)-entries, and is also closed under
multiplication. Both ideals are proper, so both are prime.
If \(H\subsetneq B\), choose \(z\notin H\).
Then \(X_0,Z_z\notin M_H\) but \(X_0Z_z=Z_0\in M_H\);
thus \(M_H\) is not prime.

Suppose \(I_{J,H}\) is prime and \(H\subsetneq B\).
The rule \(Z_zZ_w=Z_{zw}\) shows that \(H\) is prime in \(B\).
The condition \(JB\subseteq H\) gives \(J\subseteq H\cap A\).
If \(a\in H\cap A\), then \(X_aZ_1=Z_a\in I_{J,H}\)
while \(Z_1\notin I_{J,H}\), so primality gives \(a\in J\).
Thus \(I_{J,H}=Q_H\).
Conversely take a prime \(\mathfrak q\) in \(B\), and let
\(\mathfrak p=\mathfrak q\cap A\). This contraction is proper
and prime: \(1\notin\mathfrak q\), and \(ab\in\mathfrak q\)
forces \(a\in\mathfrak q\) or \(b\in\mathfrak q\).
Two \(X\)- or \(Y\)-entries outside \(Q_{\mathfrak q}\)
have coefficients outside \(\mathfrak p\), so their product
coefficient, with either sign, remains outside \(\mathfrak p\).
An \(X_a\) outside this ideal and a \(Z_z\) outside it have
\(a,z\notin\mathfrak q\), so \(az\notin\mathfrak q\).
For \(Y_aZ_z\), the same holds for \(taz\), since \(t\) is a
unit and cannot lie in a proper prime. Two \(Z\)-entries outside
the ideal have product outside by primality of \(\mathfrak q\).
These are all products in its complement, proving primality.

If \(H=B\), the ideal \(I_{J,B}\) is proper precisely when
\(J\subsetneq A\). Its complement consists of \(X_a,Y_a\)
with \(a\notin J\). It is multiplicatively closed when \(J\)
is prime, by the first row of (MFC.14). Conversely, if it is
multiplicatively closed and \(ab\in J\), the product
\(X_aX_b\) belongs to the ideal, so at least one factor does,
proving that \(J\) is prime. This gives exactly the
\(T_{\mathfrak p}\)-family. Together with (MFC.21), it proves
(MFC.24).

The complete prime inclusion rules are
\[
 \begin{aligned}
 &P_\tau\subsetneq M,\quad
 P_\tau\subsetneq Q_{\mathfrak q},\quad
 P_\tau\subsetneq T_{\mathfrak p},\quad
 M\subsetneq T_{\mathfrak p},\\
 &Q_{\mathfrak q}\subseteq Q_{\mathfrak q'}
       \iff\mathfrak q\subseteq\mathfrak q',\qquad
 T_{\mathfrak p}\subseteq T_{\mathfrak p'}
       \iff\mathfrak p\subseteq\mathfrak p',\\
 &Q_{\mathfrak q}\subseteq T_{\mathfrak p}
       \iff\mathfrak q\cap A\subseteq\mathfrak p .
 \end{aligned}                                           \tag{MFC.25}
\]
Every inclusion of a \(Q\)-prime in a \(T\)-prime is strict,
because \(Z_1\) lies only in the latter.
The prime \(M\) is incomparable with each \(Q_{\mathfrak q}\):
\(Z_1\in M\setminus Q_{\mathfrak q}\) and
\(X_0\in Q_{\mathfrak q}\setminus M\).
No \(T\)-prime is contained in a \(Q\)-prime, because it contains
\(Z_1\); no \(T\)-prime is contained in \(M\), because it contains
\(X_0\). These assertions and (MFC.25) follow directly from
(MFC.23), and exhaust the possible pairs of families.
The ordinary ideal convention matters here. Each \(M_H\)
is non-subtractive: \(Z_0\in M_H\) and
\(Z_0+X_0=Z_0\in M_H\), whereas \(X_0\notin M_H\).
Each \(T_{\mathfrak p}\) is likewise non-subtractive:
\(Z_0,Z_0+X_1=Z_1\in T_{\mathfrak p}\), whereas
\(X_1\notin T_{\mathfrak p}\) because \(\mathfrak p\) is proper.
Thus restricting to subtractive primes would remove
original prime branches occurring in (MFC.24).

\paragraph{MFC.10. The dimension calculation, including incomparability.}
If \(A\ne0\) and its Krull dimension is \(\dim A=d<\infty\), then
\[
                     \dim D_A=d+2.                       \tag{MFC.26}
\]
Here dimension means the ideal-theoretic prime-chain dimension
of MFC.1. We give the integral-extension argument needed for
the upper bound in full.
The ring \(B=A[t]/(t^2+1)\) is integral over its embedded copy
of \(A\), since every \(b=a+tc\) satisfies
\[
             b^2-2ab+(a^2+c^2)=0.                        \tag{MFC.27}
\]
Suppose primes \(\mathfrak q\subsetneq\mathfrak q'\) of \(B\)
had the same contraction \(\mathfrak p\) in \(A\).
Then \(R=A/\mathfrak p\) embeds in the domain
\(C=B/\mathfrak q\), and \(C\) is integral over \(R\):
reduce any monic equation over \(A\) modulo these ideals.
The nonzero prime \(\mathfrak r=\mathfrak q'/\mathfrak q\)
of \(C\) has zero contraction in \(R\).
Let \(S=R\setminus\{0\}\) and \(K=S^{-1}R\).
The localization \(S^{-1}C\) is a domain, and every \(c/s\)
in it is integral over the field \(K\). Indeed, dividing a
monic equation for \(c\) by the corresponding power of \(s\)
produces a monic equation for \(c/s\) with coefficients in \(K\).
A domain integral over a field is a field: for a nonzero
element \(v\), choose a monic relation of least degree over
that field. Its constant coefficient is nonzero, because
otherwise division by \(v\) in the domain would give a monic
relation of lower degree. Rearranging the relation expresses
\(v^{-1}\) as a polynomial in \(v\) divided by that nonzero
constant coefficient. Thus \(S^{-1}C\) is a field.
Since \(\mathfrak r\cap S=\varnothing\), its localization is
a proper prime ideal of this field: if \(1=c/s\) with
\(c\in\mathfrak r,s\in S\), clearing the equality gives
\(v(c-s)=0\) for some \(v\in S\), and the domain property
implies \(s=c\in\mathfrak r\), a contradiction.
It follows that \(S^{-1}\mathfrak r=0\). But a nonzero
\(c\in\mathfrak r\) remains nonzero after localization,
since \(sc=0\), \(s\in S\), is impossible in the domain.
This contradiction proves that strict prime inclusions in
\(B\) have strict contractions in \(A\).

The ring \(A\ne0\) has a maximal proper ideal: the union of
a chain of proper ideals is proper, since otherwise it
contains \(1\) at some stage, and Zorn's lemma applies.
A maximal ideal is prime because its quotient is a field:
if the class of \(a\) is nonzero, maximality of the ideal
generated by that class supplies its multiplicative inverse.
Thus \(d\ge0\), and a chain of length \(d\) exists. Indeed the
supremum of a nonempty set of nonnegative integer lengths,
when finite, is attained. Choose
\(\mathfrak p_0\subsetneq\cdots\subsetneq\mathfrak p_d\).
Then (MFC.25) gives the strict chain
\[
 P_\tau\subsetneq M\subsetneq T_{\mathfrak p_0}
       \subsetneq\cdots\subsetneq T_{\mathfrak p_d},        \tag{MFC.28}
\]
of length \(d+2\).

For the upper bound, it suffices to bound every finite strict
chain, since an infinite chain would contain finite chains
of arbitrarily large length. Adjoin \(P_\tau\) at its beginning
if it is missing. By (MFC.25), a chain containing \(M\) contains
no \(Q\)-prime, and after \(M\) contains only a \(T\)-chain.
If that \(T\)-chain has \(\ell\) entries, then
\(\ell-1\le d\), and the full chain has length at most
\(\ell+1\le d+2\). A chain with no \(Q\)- or \(T\)-entries
has length at most one.
A chain with \(Q\)-entries and no \(M\) has a \(Q\)-chain
of \(k\ge1\) entries followed by a \(T\)-chain of
\(\ell\ge0\) entries. The contractions of the \(Q\)-entries
are strictly increasing by the preceding proof. If
\(\ell\ge1\), they are followed by the strictly increasing
\(T\)-indices, with at most one equality, at the junction.
Therefore this chain in \(\operatorname{Spec}A\) has at least
\(k+\ell-1\) distinct entries, so \(k+\ell-2\le d\).
The original chain, including \(P_\tau\), has length
\(k+\ell\le d+2\). If \(\ell=0\), then \(k-1\le d\) and
its length is \(k\le d+1\). A chain having \(T\)-entries
but neither \(Q\)-entries nor \(M\) has length at most \(d+1\).
These cases exhaust (MFC.24), proving the upper bound and
therefore (MFC.26).

\paragraph{MFC.11. Higher distributive-lattice zero fibres.}
Let \(L\) be a bounded distributive lattice, with bottom \(0_L\)
and top \(1_L\), and let \(R\ne0\) be a ring. The original carrier is
\[
 \begin{split}
 G_L(R)&=\{(0_R,\lambda):\lambda\in L\}
                \cup\{(r,1_L):r\in R\}\subseteq R\times L,\\
 (r,\lambda)+(s,\mu)&=(r+s,\lambda\vee\mu),\\
 (r,\lambda)(s,\mu)&=(rs,\lambda\wedge\mu),\\
 0_{G_L(R)}&=(0_R,0_L),\qquad1_{G_L(R)}=(1_R,1_L).
 \end{split}                                             \tag{MFC.29}
\]
This is a union inside a product; the two descriptions of
\((0_R,1_L)\) denote one element. A nonzero amplitude is
permitted only at top support.
The operations are closed. If \(r+s\ne0\), at least one
summand has nonzero amplitude and hence top support, so the
sum has top support. If \(rs\ne0\), both amplitudes are nonzero
and both supports are top. Zero-amplitude results allow any
support. Associativity and commutativity follow in each
coordinate from the ring and lattice laws. Distributivity
follows from ring distributivity and distributivity of meet
over join in \(L\). The displayed zero and unit satisfy their
laws in each coordinate. Thus (MFC.29) is a semiring.

The projections
\[
 p_L:G_L(R)\to R,\quad(r,\lambda)\mapsto r,\qquad
 \chi_L:G_L(R)\to L,\quad(r,\lambda)\mapsto\lambda,
 \quad p_L^{-1}(0_R)=\{(0_R,\lambda):\lambda\in L\}         \tag{MFC.30}
\]
are unital semiring homomorphisms, where \(L\) has join
and meet as its operations. The last fibre is exactly
the entire original lattice.
Moreover \(p_L\) is the universal unital homomorphism
from \(G_L(R)\) to a ring. To prove this, write
\(z_\lambda=(0_R,\lambda)\) and \(e=(0_R,1_L)\).
For any homomorphism \(f:G_L(R)\to T\) to a ring,
\(e+z_\lambda=e\) gives \(f(z_\lambda)=0_T\) by
additive cancellation in \(T\).
The function \(\bar f:R\to T\), \(r\mapsto f(r,1_L)\),
therefore preserves zero, addition, multiplication, and unit,
so is a ring homomorphism. Every element of \(G_L(R)\)
is a \(z_\lambda\) or a \((r,1_L)\), proving
\(f=\bar f p_L\), and the surjectivity of \(p_L\) proves
uniqueness. If \(L\) is the one-element lattice, the carrier
is \(R\) itself; the distinct absent and supported-zero
elements occur precisely when \(0_L\ne1_L\).

\paragraph{MFC.12. All lattice-localization fibres and the universal map.}
For every \(\lambda\in L\), the element \(z_\lambda\)
is idempotent and
\[
 z_\lambda G_L(R)=\{z_\nu:\nu\le\lambda\},\qquad
 G_L(R)[z_\lambda^{-1}]\cong\downarrow\lambda,\qquad
 \ell_\lambda:G_L(R)\to\downarrow\lambda,\quad
 (r,\mu)\mapsto\lambda\wedge\mu.                           \tag{MFC.31}
\]
Here \(\downarrow\lambda\) has join and meet, bottom \(0_L\),
and unit \(\lambda\). Multiplying in (MFC.29) gives
\(z_\lambda(r,\mu)=z_{\lambda\wedge\mu}\);
every \(\nu\le\lambda\) occurs by taking \((0_R,\nu)\).
This proves the corner equality and its identification
with \(\downarrow\lambda\), including its unit.
Distributivity of \(L\) proves that \(\ell_\lambda\)
preserves addition, and associativity and idempotence of
meet prove that it preserves multiplication. It sends
the original zero to \(0_L\), the original unit to
\(\lambda\), and \(z_\lambda\) to that unit.
For a unital map \(f\) making \(z_\lambda\) invertible,
idempotence forces \(f(z_\lambda)=1\). Hence
\(f(x)=f(z_\lambda x)\), and restriction to the corner
gives its unique factorization through \(\ell_\lambda\).
This is the required localization universal property.
At \(\lambda=0_L\) the localization is the one-element
semiring; the universal statement allows that zero semiring.
Its exact fibres are
\[
 \ell_\lambda^{-1}(\nu)
 =\{(r,\mu)\in G_L(R):\lambda\wedge\mu=\nu\},
 \qquad \nu\le\lambda.                                   \tag{MFC.32}
\]
Thus every top-supported amplitude, including every nonzero
amplitude, maps to the unit \(\lambda\) of the localized
interval. The formula specifies every erased amplitude
and every identified support stratum.

\paragraph{MFC.13. The higher-zero carrier inside the full independent cube.}
Take \(R=R_I\) and \(L=\mathcal P(I)\). Define
\[
 J:G_{\mathcal P(I)}(R_I)\hookrightarrow P_I,\qquad
 J(r,I)=(r_i^\bullet)_i,\qquad
 J(0,A)=x(0,A).                                          \tag{MFC.33}
\]
The two formulas agree at \((0,I)\).
Two full-support entries add and multiply by the ring
operations in every coordinate, so \(J\) preserves their
sum and product. Two zero-amplitude entries have support
union under addition and intersection under multiplication,
by (MFC.5), again as in (MFC.29).
A full entry plus a zero-amplitude entry stays full with
the original amplitude; a full entry multiplied by a
zero-amplitude entry gives the same zero-amplitude face,
because supported multiplication by zero gives supported
zero in exactly the active coordinates of that face.
These cases exhaust all pairs and prove preservation
of addition and multiplication. The bottom entry maps
to the all-absent zero and \((1,I)\) to the unit.
Recovering amplitude and support proves injectivity.
Its exact image is all full-support tuples together
with all zero-amplitude faces. Proper faces with nonzero
amplitude are the additional elements of \(P_I\).
The synchronized inclusion factors as
\[
 G(R_I)\xrightarrow{\ k\ }G_{\mathcal P(I)}(R_I)
                 \xrightarrow{\ J\ }P_I,\qquad
 k(\tau_{R_I})=(0,\varnothing),\quad
 k(r^\bullet)=(r,I),\quad Jk=j.                           \tag{MFC.34}
\]
The formulas prove that \(k\) is injective and unital,
using (MFC.2) and (MFC.29). Both maps also commute
with amplitude and support, since those coordinates
are the coordinates used in their definitions.

\paragraph{MFC.14. Every original chain stratum and each threshold fibre.}
Let \(n\ge1\) and \(C_n=\{0<1<\cdots<n\}\).
In \(G_{C_n}(R)\) put
\[
 \begin{gathered}
 e=(0_R,n),\qquad \tau_i=(0_R,n-i)\quad(1\le i\le n),\\
 \tau_0=e\quad\text{when indexing formulas},\\
 \tau_i+\tau_j=\tau_{\min(i,j)},\qquad
 \tau_i\tau_j=\tau_{\max(i,j)} .
 \end{gathered}                                          \tag{MFC.35}
\]
These equalities follow from join \(=\max\) and meet \(=\min\)
in (MFC.29). For every \(r\in R\),
\((r,n)+\tau_i=(r,n)\) and \((r,n)\tau_i=\tau_i\).
The carrier consists of the full ring copy
\(\{(r,n):r\in R\}\), with its zero \(e\), and the
\(n\) further distinct entries \(\tau_1,\ldots,\tau_n\).
In particular its amplitude-zero fibre has \(n+1\) entries.

For \(1\le k\le n\), the original threshold and induced map are
\[
 \begin{gathered}
 \alpha_k:C_n\to\mathbb B,\qquad
 \alpha_k(j)=
 \begin{cases}0&j<k,\\1&j\ge k,\end{cases}\\
 T_k:G_{C_n}(R)\to G(R),\qquad T_k(r,n)=r^\bullet,\qquad
 T_k(0_R,j)=
 \begin{cases}\tau_R&j<k,\\e_R&j\ge k.\end{cases}
 \end{gathered}                                          \tag{MFC.36}
\]
The formulas agree at \((0_R,n)\).
The threshold sends bottom to zero and top to one.
The maximum of two indices is at least \(k\) precisely
when at least one is, and the minimum is at least \(k\)
precisely when both are. Thus it preserves join and meet.
For two zero-amplitude entries this proves both laws
for \(T_k\). Two full entries use the original ring
operations and stay full, even upon amplitude cancellation.
A full entry plus any zero-amplitude entry stays full;
its image plus either \(\tau_R\) or \(e_R\) is the same
supported image. A full entry times \((0_R,j)\) equals
\((0_R,j)\); its supported image times the prescribed
image of that entry is respectively \(\tau_R\) or \(e_R\).
These are all cases, proving that \(T_k\) is a unital
semiring homomorphism.
Its entire fibres are
\[
 \begin{aligned}
 T_k^{-1}(\tau_R)&=\{(0_R,j):0\le j<k\},\\
 T_k^{-1}(e_R)&=\{(0_R,j):k\le j\le n\},\\
 T_k^{-1}(r^\bullet)&=\{(r,n)\}\quad(r\ne0_R).
 \end{aligned}                                           \tag{MFC.37}
\]
They have respectively \(k,n-k+1,1\) entries. These formulas
retain the exact depth lost under every binary threshold.

\paragraph{MFC.15. The all-ideal chain for the original two arithmetic rings.}
Let \(p\) be a prime integer, let \(k\) now denote a field,
and let \(n\ge1\). Define the two rings and their ideals by
\[
 R_{p,n}=\mathbb Z/p^n\mathbb Z,\quad
 I_a=p^aR_{p,n},\qquad
 R_{k,n}=k[\epsilon]/(\epsilon^n),\quad
 J_a=\epsilon^aR_{k,n},\qquad 0\le a\le n.                \tag{MFC.38}
\]
In particular \(I_0=R_{p,n}\), \(I_n=(0)\),
\(J_0=R_{k,n}\), and \(J_n=(0)\).
We prove that these are all ideals, not a selected subfamily.

For a nonzero element of \(R_{p,n}\), take its representative
between \(1\) and \(p^n-1\). Its exact power of \(p\)
is \(p^a\) with \(0\le a<n\); thus the element is \(p^a u\)
with \(u\) prime to \(p\). The class of \(u\) is a unit:
the Euclidean algorithm applied to \(u,p^n\), whose greatest
common divisor is \(1\), supplies integers \(v,w\) with
\(uv+p^nw=1\). Let \(\mathfrak a\) be a nonzero ideal,
and choose among its nonzero elements one with least
exponent \(a\). Multiplying this element by \(u^{-1}\)
puts \(p^a\) in \(\mathfrak a\), so \(I_a\subseteq\mathfrak a\).
Every other nonzero element has exponent at least \(a\)
and is therefore a multiple of \(p^a\); zero also lies
in \(I_a\). Hence \(\mathfrak a=I_a\). The zero ideal
is \(I_n\). The ideals are distinct: for \(a<n\),
an equality \(p^a=p^{a+1}b\) modulo \(p^n\) would imply
that \(p^{n-a}\) divides \(1-pb\), impossible upon
reduction modulo \(p\).

Every element of \(R_{k,n}\) has a unique polynomial
representative of degree less than \(n\), by division
by the monic polynomial \(\epsilon^n\). A nonzero
representative with least nonzero degree \(a\) is
\(\epsilon^a u\), where \(u=u_0(1+v)\), \(u_0\in k^\times\),
and \(v\in(\epsilon)\). Since \(v^n=0\),
\[
        u^{-1}=u_0^{-1}\sum_{j=0}^{n-1}(-v)^j             \tag{MFC.39}
\]
is an inverse: multiplication by \(1+v\) gives
\(1-(-v)^n=1\).
In a nonzero ideal choose an element with least such
degree \(a\). Multiplication by \(u^{-1}\) puts
\(\epsilon^a\) in the ideal, while every other element
has degree at least \(a\) and is divisible by
\(\epsilon^a\). Thus the ideal is \(J_a\); the zero
ideal is \(J_n\). Uniqueness of the degree-less-than-\(n\)
representative shows that
\(\epsilon^a\notin J_{a+1}\) for \(a<n\), proving that
these \(n+1\) ideals are distinct.

In both rings the ideals are nested by reversed exponent.
Consequently their sums and intersections are
\[
 I_a+I_b=I_{\min(a,b)},\quad I_a\cap I_b=I_{\max(a,b)},
 \qquad
 J_a+J_b=J_{\min(a,b)},\quad J_a\cap J_b=J_{\max(a,b)}.
                                                               \tag{MFC.40}
\]
For instance, if \(a\le b\), then \(I_b\subseteq I_a\),
so the sum is \(I_a\) and the intersection is \(I_b\);
the other cases are the same two inclusions with indices
interchanged. It follows that the explicit bijections
\[
 C_n\xrightarrow{\ \theta_p\ }\operatorname{Id}(R_{p,n}),
       \quad j\mapsto I_{n-j},\qquad
 C_n\xrightarrow{\ \theta_k\ }\operatorname{Id}(R_{k,n}),
       \quad j\mapsto J_{n-j}                             \tag{MFC.41}
\]
preserve bottom, top, join, and meet, with ideal sum as
join and ideal intersection as meet. Surjectivity here is
the all-ideal statement just proved; injectivity is the
strictness proved above. Thus these are isomorphisms of
the bounded-lattice semirings, preserving every stratum.

\paragraph{MFC.16. The separate, complete ideal-product law.}
The very same ideal carriers also retain ordinary ideal
multiplication, with its different operation:
\[
 \begin{gathered}
 I_aI_b=I_{\min(n,a+b)},\qquad
 J_aJ_b=J_{\min(n,a+b)},\\
 V_n=\bigl(\{0,\ldots,n\},\ \min,\ (a,b)\mapsto
                   \min(n,a+b),\ 0_{V_n}=n,\ 1_{V_n}=0\bigr).
 \end{gathered}                                          \tag{MFC.42}
\]
Indeed the product of the principal ideals generated by
\(x,y\) is generated by \(xy\): every finite sum of
products \(xr\cdot ys\) is a multiple of \(xy\), and
\(xy\) itself is one of the generating products.
Applying this fact to \(p^a,p^b\), and to
\(\epsilon^a,\epsilon^b\), proves the displayed laws,
including the vanishing for \(a+b\ge n\).
The maps \(a\mapsto I_a\) and \(a\mapsto J_a\) are
therefore bijective semiring homomorphisms from \(V_n\)
to the respective all-ideal carriers with sum and
ordinary ideal product. For completeness, associativity
of the operation on exponents follows from
\(\min(n,\min(n,a+b)+c)=\min(n,a+b+c)\);
its unit and absorbing zero are \(0\) and \(n\).
The identity
\[
 \min(n,a+\min(b,c))
       =\min\bigl(\min(n,a+b),\,\min(n,a+c)\bigr)
                                                               \tag{MFC.43}
\]
proves distributivity over the addition \(\min\).
This also checks all semiring laws directly on \(V_n\).

Under the order-preserving coordinates of (MFC.41),
ordinary ideal multiplication is the operation
\[
             j\mathbin{\odot}\ell=\max(0,j+\ell-n)
             \quad\text{on }C_n,                         \tag{MFC.44}
\]
because \(n-\min(n,(n-j)+(n-\ell))=\max(0,j+\ell-n)\).
It is distinct from the meet \(\min(j,\ell)\) when
\(n\ge2\): for \(j=\ell=n-1\), their values are
\(\max(0,n-2)\) and \(n-1\). Equivalently,
\(I_1^2=I_2\) whereas \(I_1\cap I_1=I_1\), with
the convention \(I_2=0\) when \(n=2\), and the
same distinction holds for \(J_1\).
The identity map on either all-ideal carrier preserves
the carrier, order, zero, unit, and addition between
these two presentations, but does not in general
preserve multiplication. Equations (MFC.40)--(MFC.44)
give its exact transported operations. No identification
of the two underlying arithmetic rings is asserted.

\paragraph{MFC.17. Coefficient transport, support preservation, and prime maps.}
The preceding constructions are accompanied by their original
coefficient maps; they are not merely lists of properties.
For a unital ring homomorphism \(f:R\to R'\), define
\[
 G(f)(\tau_R)=\tau_{R'},\qquad
 G(f)(r^\bullet)=f(r)^\bullet.                             \tag{MFC.45}
\]
Equation (MFC.2) proves that this is a unital semiring
homomorphism, including when a nonzero \(r\) has
\(f(r)=0\): that entry is still supported.
For a family \(f_i:R_i\to R_i'\), the product map
\(\prod_iG(f_i)\) preserves each mask exactly, has
amplitude \(\prod_i f_i\), and commutes with \(j\)
and \(J\) in (MFC.8) and (MFC.33).
Each assertion follows by substituting into one
coordinate of the displayed formulas; inactive
coordinates stay absent and active ones stay supported.
Thus the support square and every independent face
remain part of the transported diagram.

For \(f:A\to A'\), let
\(\bar f:A[t]/(t^2+1)\to A'[t]/(t^2+1)\) be
\(\bar f(a+tb)=f(a)+tf(b)\). This is well-defined
and a unital ring homomorphism: applying \(f\) to
the two coefficients of a product
\((ac-bd)+t(ad+bc)\) gives exactly the corresponding
product of the two images. Define
\[
 D(f)(a,b)=(G(f)(a),G(f)(b)),\quad
 q_{D_{A'}}D(f)=G(\bar f)q_{D_A},\quad
 D(f)\iota_A=\iota_{A'}G(\bar f).                          \tag{MFC.46}
\]
The equality \(G(f)(u)=(-1_{A'})^\bullet\) and
(MFC.11) prove that \(D(f)\) is a unital semiring
homomorphism. It sends \(X_a,Y_b,Z_z\) respectively
to \(X_{f(a)},Y_{f(b)},Z_{\bar f(z)}\), so it
preserves each of the four faces exactly and proves
both commuting identities in (MFC.46). It sends
\(Z_1\) to \(Z_1\), hence also induces the corner
and localization map already represented by \(G(\bar f)\).

For ideals in the primed rings its inverse-image formulas are
\[
 \begin{aligned}
 D(f)^{-1}(M_{H'})&=M_{\bar f^{-1}(H')},\\
 D(f)^{-1}(I_{J',H'})&=I_{f^{-1}(J'),\,\bar f^{-1}(H')},\\
 D(f)^{-1}(P_\tau)&=P_\tau,\qquad D(f)^{-1}(M)=M,\\
 D(f)^{-1}(Q_{\mathfrak q'})
      &=Q_{\bar f^{-1}(\mathfrak q')},\qquad
 D(f)^{-1}(T_{\mathfrak p'})
      =T_{f^{-1}(\mathfrak p')}.
 \end{aligned}                                           \tag{MFC.47}
\]
The first two formulas follow by the exact face images,
and the required condition on the second ideal pair holds:
if \(a\in f^{-1}(J')\) and \(b\in B\), then
\(\bar f(ab)=f(a)\bar f(b)\in J'B'\subseteq H'\).
Closure under sums gives
\(f^{-1}(J')B\subseteq\bar f^{-1}(H')\).
The two constant-family formulas follow because nonempty
faces and full faces remain in precisely those classes.
For the \(Q\)-formula the contraction identity
\(\bar f^{-1}(\mathfrak q')\cap A
=f^{-1}(\mathfrak q'\cap A')\) proves agreement of both
indices. Inverse images of proper prime ideals under
unital ring maps are proper primes, because they omit
the unit and satisfy the prime product implication.
This proves the typing and the final two formulas.
These maps are continuous on ideal-theoretic spectra:
for any ideal \(\mathcal I\subseteq D_A\), the inverse
image, under prime contraction, of the closed set of
primes containing \(\mathcal I\) is the set of primes
of \(D_{A'}\) containing the ideal generated by
\(D(f)(\mathcal I)\). The two containment conditions
are equivalent by the definition of inverse image.

Finally, for a bounded-lattice homomorphism
\(\beta:L\to L'\) preserving bottom, top, joins, and meets and a
unital ring homomorphism \(f:R\to R'\), the map
\[
 G_\beta(f):G_L(R)\to G_{L'}(R'),\quad
 (r,\lambda)\mapsto(f(r),\beta(\lambda))                  \tag{MFC.48}
\]
is well-defined: a nonzero target amplitude comes
from a nonzero source amplitude, hence from
\(\lambda=1_L\), whose image is \(1_{L'}\).
The two operations in (MFC.29) and the hypotheses
on \(f,\beta\) prove preservation of sum, product,
zero, and unit. It commutes with amplitude and
support via \(f,\beta\), and sends \(z_\lambda\)
to \(z_{\beta(\lambda)}\).
Its localization transport is the map
\(\downarrow\lambda\to\downarrow\beta(\lambda)\),
\(\nu\mapsto\beta(\nu)\), because
\[
 \beta\bigl(\ell_\lambda(r,\mu)\bigr)
 =\beta(\lambda\wedge\mu)
 =\beta(\lambda)\wedge\beta(\mu)
 =\ell_{\beta(\lambda)}
       \bigl(G_\beta(f)(r,\mu)\bigr).                      \tag{MFC.49}
\]
It preserves the interval unit and bottom.
The threshold maps of MFC.14 are exactly this
construction for \(f=\operatorname{id}_R\) and
\(\beta=\alpha_k\), after the isomorphism
\(G_{\mathbb B}(R)\cong G(R)\) sending
\((0_R,0)\) to \(\tau_R\) and \((r,1)\) to
\(r^\bullet\). Substitution in (MFC.29) and (MFC.2)
proves that this last bijection is a unital
semiring isomorphism.
All maps in this attachment preserve their stated
source indices and codomains. Their exact fibres
identify the information forgotten by a specified
observation while retaining the full carriers to
which the observation is attached.
